\documentclass[10pt]{article}
\usepackage[margin=0.76in]{geometry}
\usepackage{amsmath,amssymb,amsthm,mathtools}
\usepackage{microtype}
\usepackage[hidelinks]{hyperref}
\usepackage{enumitem}

\newtheorem{theorem}{Theorem}
\newtheorem{lemma}[theorem]{Lemma}
\newtheorem{remark}[theorem]{Remark}
\newcommand{\C}{\mathbb C}
\newcommand{\N}{\mathbb N}
\newcommand{\calF}{\mathcal F}
\newcommand{\calH}{\mathcal H}
\newcommand{\tr}{\operatorname{tr}}
\newcommand{\Ran}{\operatorname{Ran}}
\newcommand{\Id}{\operatorname{Id}}
\newcommand{\diag}{\operatorname{diag}}
\newcommand{\Span}{\operatorname{span}}
\newcommand{\supp}{\operatorname{supp}}

\begin{document}
\begin{center}
{\Large\bfseries Sharp sampling numbers for nonseparable\par
reproducing-kernel Hilbert spaces\par}
\medskip
Automated Math Discovery Run \texttt{fa\_banach\_001}; Model: GPT5.6\par
9 August 2026
\end{center}

\begin{center}
\fbox{\begin{minipage}{0.92\textwidth}
\textbf{Status: full proof candidate.}
Remark~13 of Dolbeault--Krieg--Ullrich predicts that separability can be
removed from their sharp $L_2$ sampling theorem by adding
$\tr_0(K)/m$.  The theorem below proves exactly that prediction.  The
proof is nonconstructive and the displayed universal constant is intentionally
not optimized.
\end{minipage}}
\end{center}

\section*{1. Source question and statement}
Let $(D,\mathcal A,\mu)$ be a measure space and let $H=H(K)$ be a (not
necessarily separable) RKHS of measurable complex-valued functions on $D$,
with $K$ measurable on $(D\times D,\mathcal A\otimes\mathcal A)$.
Assume
\[
  \tr(K):=\int_D K(x,x)\,d\mu(x)<\infty.
\]
Then the canonical map $T:H\to L_2(\mu)$ is Hilbert--Schmidt.  Write its
positive singular values, with multiplicity, as
$\sigma_0\geq\sigma_1\geq\cdots>0$, and put $\sigma_k=0$ after the end
when the rank is finite.  For the unit ball $F=B_H$, use the source convention
$d_k(F)=\sigma_k$ for its $L_2$ Kolmogorov widths.

Let $H_0=\ker T$, let $K_0$ be its reproducing kernel, and set
\[
 \tau_0:=\tr_0(K):=\int_DK_0(x,x)\,d\mu(x)
       =\tr(K)-\sum_{k\geq0}\sigma_k^2\geq0.                 \tag{1}
\]
The sampling number $g_n(F)$ is the least worst-case $L_2$ error of a
linear algorithm using at most $n$ point values.

\begin{theorem}[nonseparable sharp sampling bound]\label{thm:main}
There is a universal integer $C$ such that, for every $n\geq1$,
\[
 g_{Cn}(B_H)^2
 \leq \frac1n\left(\sum_{k\geq n}d_k(B_H)^2+\tr_0(K)\right).       \tag{2}
\]
One admissible value furnished by the proof is
$C=56{,}116{,}800$.
\end{theorem}

For $\tau_0=0$, this is the source theorem (up to an inessential change of
its universal constant).  Printed p.~7, Remark~13, of the source explicitly
predicts (2): remove separability and put $\tr_0(K)/m$ inside the square root
on the right-hand side \cite[p.~7, Remark~13]{DKU}.

\section*{2. Idea of the proof}
The compact spectral theorem splits
\[
 H=H_1\oplus H_0,
 \qquad H_1=(\ker T)^\perp,
\]
where $H_1$ is separable and carries all nonzero singular directions, while
every $h\in H_0$ is zero in $L_2$.  It is tempting simply to discard $H_0$,
but that is wrong: a function which vanishes almost everywhere can still have
nonzero values at the finitely many sampling nodes and can corrupt a least
squares fit.

We add the normalized diagonal $K_0(x,x)/\tau_0$ as a third component of the
source's leverage-score density, alongside the leading and tail spectral
components.  This keeps the separable feature-vector norm below $3m$ and
ensures
\[
 \frac{K_0(x,x)}{\rho(x)}\leq3\tau_0.                         \tag{3}
\]
For independent nodes, the off-diagonal values $K_0(x_i,x_j)$ vanish almost
surely.  Hence the entire weighted $H_0$ sampling Gram matrix is diagonal and
has norm at most $3\tau_0$.  The source's Kadison--Singer step then reduces an
arbitrarily large good random draw to $O(m)$ nodes, while its lower frame bound
turns (3) into the desired $\tau_0/m$ error.  The spectral tail is handled
exactly as in the source.

\section*{3. A parameterized subsampling lemma}
We use the equal-weight Kadison--Singer partition lemma employed in
Section~4 of \cite{DKU}.  The following is the same argument with the norm
budget changed from $2m$ to $3m$ and with a leading block of dimension
$r\leq m$.

\begin{lemma}\label{lem:sparse}
Let $N>600m$, $1\leq r\leq m$, and
$y_1,\ldots,y_N\in\C^r\oplus\ell_2$ satisfy $\|y_i\|_2^2\leq3m$ and
\[
 \left\|\frac1N\sum_{i=1}^N y_i y_i^*
       -\begin{pmatrix}I_r&0\\0&\Lambda\end{pmatrix}\right\|
 \leq\frac12,
 \qquad \|\Lambda\|\leq1.                                  \tag{4}
\]
There is $J\subseteq\{1,\ldots,N\}$ with $|J|\leq64800m$ such that
\[
 \left(\sum_{i\in J}y_i y_i^*\right)_{<r}\geq75mI_r,
 \qquad
 \sum_{i\in J}y_i y_i^*\leq32400mI.                         \tag{5}
\]
\end{lemma}

\begin{proof}
Compress the lower coordinates to the span of the $N$ lower parts of the
$y_i$.  This makes the problem finite-dimensional without changing the
vectors or the leading block.  If
$E=\diag(I_r,\Lambda)$, split the positive matrix $I-E$ into rank-one
terms, repeat each term sufficiently many times, and scale the copies so that
the artificial vectors have squared norm at most $\delta:=3m$.  They have
zero leading coordinates.  As in \cite[Section~4]{DKU}, after adding them the
total frame operator $M$ obeys
\[
 \frac N2I\leq M\leq\frac{3N}{2}I.
\]

The partition lemma used in the source says that if
$\alpha I\leq M\leq\beta I$ and
$\alpha>100\delta$, the vectors can be partitioned into blocks $J_s$ with
\[
 25\delta I\leq\sum_{i\in J_s}y_i y_i^*
 \leq3600\frac\beta\alpha\delta I.
\]
Here $\alpha=N/2$, $\beta=3N/2$, and $\delta=3m$, so every block lies
between $75mI$ and $32400mI$.  If there are $t$ blocks, the lower bound
$M\geq(N/2)I$ and the block upper bounds give
$t\geq N/(64800m)$.  Some block therefore contains at most
$N/t\leq64800m$ of the original vectors.  Retain those original vectors and
delete its artificial vectors.  Deletion preserves the upper bound, and the
leading lower bound is unchanged because all artificial vectors have zero
leading coordinates.  Undoing the compression proves (5).
\end{proof}

\section*{4. Proof of Theorem~\ref{thm:main}}
\begin{proof}
If $T=0$, the zero algorithm has zero $L_2$ error, so assume $T\neq0$.
Choose an orthonormal basis $(\phi_k)$ of $H_1$ consisting of right singular
vectors and an orthonormal system $(b_k)$ in $L_2$ such that
\[
 \phi_k=\sigma_k b_k
\]
as functions.  This countable system gives the pointwise RKHS expansion of
the separable space $H_1$.

Fix a sampling budget parameter $m\geq1$ and put
$r=\min\{m,\operatorname{rank}T\}$.  Define the active probability-density
components
\[
 A(x)=\frac1r\sum_{k<r}|b_k(x)|^2,
 \quad
 B(x)=\frac1S\sum_{k\geq r}\sigma_k^2|b_k(x)|^2,
 \quad
 C_0(x)=\frac{K_0(x,x)}{\tau_0},                              \tag{6}
\]
where $S=\sum_{k\geq r}\sigma_k^2$, and omit $B$ if $S=0$ and $C_0$ if
$\tau_0=0$.  If $q\leq3$ components remain, let $\rho$ be their average.
Each active component integrates to one, so $d\nu=\rho\,d\mu$ is a
probability measure.  In particular,
$K_0(x,x)/\rho(x)\leq q\tau_0\leq3\tau_0$ whenever $\tau_0>0$.

If $S>0$, put
\[
 \gamma^2=\max\left\{\sigma_r^2,\frac Sm\right\};             \tag{7}
\]
if $S=0$, put $\gamma=0$ and omit all tail coordinates.  For $S>0$ define
the separable feature vector
\[
 y(x)=\rho(x)^{-1/2}
 \left((b_k(x))_{k<r},\;(\gamma^{-1}\sigma_kb_k(x))_{k\geq r}\right).
                                                                    \tag{8}
\]
Since $r\leq m$ and $S/\gamma^2\leq m$, (6) gives
\[
 \|y(x)\|_2^2
 =\frac{rA(x)+(S/\gamma^2)B(x)}{\rho(x)}\leq qm\leq3m.       \tag{9}
\]
Orthogonality of the $b_k$ gives
\[
 \mathbb E_\nu[y(X)y(X)^*]
 =\diag\left(I_r,(\sigma_k^2/\gamma^2)_{k\geq r}\right),     \tag{10}
\]
whose lower block has norm at most one.

The random trace-class operator $y(X)y(X)^*$ has trace norm at most $3m$.
The strong law in the separable trace class therefore yields an infinite iid
sequence for which the empirical covariance converges in trace norm, hence
in operator norm, to (10).  Also, for each fixed $z\in D$,
$K_0(\cdot,z)\in H_0$ and is zero $\mu$-almost everywhere.  Since
$\nu\ll\mu$, Tonelli's theorem gives
\[
 K_0(X_i,X_j)=0\quad(i\neq j)                                \tag{11}
\]
almost surely for an iid sequence.  Choose one realization satisfying (11)
and then $N>600m$ large enough that (4) holds.  Lemma~\ref{lem:sparse}
produces a subset $J$ of at most $64800m$ nodes satisfying (5).

Let
\[
 G=\left(\frac{b_k(x_i)}{\sqrt{\rho(x_i)}}\right)_{i\in J,k<r},
 \qquad
 \Phi=\left(\frac{\sigma_kb_k(x_i)}{\sqrt{\rho(x_i)}}
       \right)_{i\in J,k\geq r}.
\]
The two conclusions of (5) imply
\[
 G^*G\geq75mI_r,
 \quad \|G^+\|^2\leq\frac1{75m},
 \quad \|\Phi\|^2\leq32400m\gamma^2.                       \tag{12}
\]
Let $R_0:H_0\to\C^J$ be weighted evaluation, defined by
\[
 (R_0h)_i=\frac{h(x_i)}{\sqrt{\rho(x_i)}}.
\]
Its row Gram matrix has $(i,j)$ entry
$K_0(x_i,x_j)/\sqrt{\rho(x_i)\rho(x_j)}$.  By (3) and (11), it is diagonal
and
\[
 \|R_0\|^2\leq3\tau_0.                                      \tag{13}
\]

Use the weighted least-squares algorithm
\[
 \mathcal A f
 =\sum_{k<r}\left(G^+
   \left(\frac{f(x_i)}{\sqrt{\rho(x_i)}}\right)_{i\in J}
   \right)_k b_k.                                             \tag{14}
\]
It is linear and exact on $\Span\{b_k:k<r\}$.  Decompose
$f=P_rf+Q_rf+f_0$ according to the leading singular directions, the
spectral tail, and $H_0$.  The residual sampling map on
$H_{1,\geq r}\oplus H_0$ is the row operator $[\Phi\;R_0]$, so (12)--(13)
give
\[
 \|[\Phi\;R_0]\|^2
 =\|\Phi\Phi^*+R_0R_0^*\|
 \leq32400m\gamma^2+3\tau_0.                                \tag{15}
\]
Pythagoras in $L_2$, exactness on the leading space, and
$\|Q_rf\|_{L_2}\leq\sigma_r\|Q_rf\|_H$ now yield, for
$\|f\|_H\leq1$,
\begin{align*}
 \|f-\mathcal Af\|_{L_2}^2
 &\leq \gamma^2
 +\frac{32400m\gamma^2+3\tau_0}{75m}\\
 &=433\gamma^2+\frac{\tau_0}{25m}.                            \tag{16}
\end{align*}
The same computation covers finite rank when the tail is absent.

Monotonicity of the singular values gives
\[
 \gamma^2\leq\frac2m
     \sum_{k\geq\lceil m/2\rceil}\sigma_k^2.                 \tag{17}
\]
Consequently,
\[
 g_{64800m}(F)^2
 \leq\frac{866}{m}\sum_{k\geq\lceil m/2\rceil}d_k(F)^2
      +\frac{\tau_0}{25m}.                                   \tag{18}
\]
Apply (18) with $m=866n$.  Since the tail beginning at $433n$ is no
larger than the tail beginning at $n$, and sampling numbers decrease with
the number of nodes,
\[
 g_{(64800)(866)n}(F)^2
 \leq\frac1n\sum_{k\geq n}d_k(F)^2+\frac{\tau_0}{n}.
\]
As $(64800)(866)=56{,}116{,}800$, this proves the theorem.
\end{proof}

\section*{5. Checks, dependence, and scope}
\begin{itemize}[leftmargin=1.5em]
\item \textbf{Null-space necessity.}  Although $H_0$ is invisible in $L_2$,
its point values can perturb least squares.  Bound (13), not deletion of
$H_0$, is what justifies the extension.
\item \textbf{Degenerate cases.}  If $T=0$, the zero algorithm is exact in
$L_2$.  If $T$ has finite positive rank, the same proof takes all positive
singular directions once $m$ exceeds the rank; (13) still gives the
$\tau_0/m$ term.  If $\tau_0=0$, the third density component and (13) vanish.
\item \textbf{External dependency.}  The only deep imported step is the
equal-weight partition lemma used by Dolbeault--Krieg--Ullrich, ultimately
coming from the Marcus--Spielman--Srivastava solution of Kadison--Singer.
The compact spectral decomposition and definition (1) are the framework of
Moeller--Ullrich \cite[Section~7]{MU}.
\item \textbf{Nonconstructive character.}  The trace-class strong law and
the Kadison--Singer partition supply existence, not an efficient algorithm.
The numerical constant is a bookkeeping constant, not a sharp claim.
\item \textbf{Scope.}  This proves only the explicit nonseparable-RKHS
extension in Remark~13.  It does not address the paper's broad statement that
$L_p$ approximation for $p\neq2$ is widely open.
\end{itemize}

\section*{6. Provenance and literature search}
Moeller--Ullrich already prove a nonseparable estimate with the same
$\tr_0(K)$ obstruction, but their Theorem~1.2 retains logarithmic
oversampling/tail loss \cite[Theorem~1.2 and Section~7]{MU}.  The present
argument combines their null-kernel density and diagonal-Gram observation
with the later sharp subsampling mechanism of \cite{DKU}; neither cited
paper states (2).

Exact-title, exact-symbol (\texttt{tr\_0(K)}), nonseparable-RKHS, and citing
arXiv searches were run on 9 August 2026.  They found no later paper stating
this theorem.  That is evidence against an obvious duplicate, not a certified
priority or exhaustive-literature claim.

\begin{thebibliography}{9}
\bibitem{DKU}
M. Dolbeault, D. Krieg, and M. Ullrich,
\emph{A sharp upper bound for sampling numbers in $L_2$},
Appl. Comput. Harmon. Anal. 63 (2023), 113--134;
\href{https://arxiv.org/abs/2204.12621}{arXiv:2204.12621}.

\bibitem{MU}
M. Moeller and T. Ullrich,
\emph{$L_2$-norm sampling discretization and recovery of functions from
RKHS with finite trace},
Sampl. Theory Signal Process. Data Anal. 19 (2021), Article 13;
\href{https://arxiv.org/abs/2009.11940}{arXiv:2009.11940}.

\bibitem{MSS}
A. W. Marcus, D. A. Spielman, and N. Srivastava,
\emph{Interlacing families II: mixed characteristic polynomials and the
Kadison--Singer problem},
Ann. of Math. 182 (2015), 327--350.
\end{thebibliography}
\end{document}
